
\documentclass[aspectratio=169]{beamer}

\mode<presentation>
\usetheme{Boadilla}
\definecolor{NTHU1}{RGB}{127,16,132}
\definecolor{NTHU2}{RGB}{228,210,231}
\definecolor{blue}{RGB}{30,90,205}
\definecolor{red}{RGB}{213,94,0}
\definecolor{green}{RGB}{0,128,0}
\definecolor{charcoalgray}{RGB}{108,104,104}
\setbeamercolor{title}{fg=NTHU1}
\setbeamercolor{frametitle}{fg=NTHU1}
\setbeamercolor{block title}{bg=NTHU1, fg=white}
\setbeamercolor{block body}{bg=NTHU2}
\setbeamercolor{structure}{fg=NTHU1}
\setbeamercolor{item projected}{fg=white}
\setbeamercolor{item}{fg=NTHU1}
\setbeamercolor{subitem}{fg=NTHU1}
\setbeamercolor{section in toc}{fg=NTHU1}
\setbeamercolor{description item}{fg=NTHU1}
\setbeamercolor{caption name}{fg=NTHU1}
\setbeamercolor{button}{bg=NTHU1, fg=white}
\setbeamercolor{caption name}{fg=NTHU1}
\usepackage{graphics}
\usepackage{tikz}
\usepackage{amsmath}
\usepackage{bbm}
\usetikzlibrary{decorations.pathreplacing}
\usepackage{geometry}
\usepackage{booktabs}
\usepackage{bookmark}
\usepackage{multirow, makecell}
\usepackage{float}
\usepackage{fancyvrb}
\usepackage{caption}
\captionsetup[figure]{singlelinecheck = false, format= hang, justification=centering}
\captionsetup[subfigure]{singlelinecheck = false, format= hang, justification=raggedright}
\usepackage{subcaption}
\usepackage{adjustbox}
\usepackage{hyperref}
\usepackage{threeparttable}
\usepackage{appendixnumberbeamer} 
\usepackage[T1]{fontenc}
\usepackage[default]{lato} 
\usepackage{smartdiagram}
\smartdiagramset{
  back arrow disabled = true,
  module minimum width=1cm,
  module x sep = 3,
  uniform arrow color=true,
}
\usepackage{xcolor}
\usepackage{colortbl}
  \definecolor{light-gray}{gray}{0.99}

\newenvironment{wideitemize}{\itemize\addtolength{\itemsep}{10pt}}{\enditemize}
\newenvironment{wideenumerate}{\enumerate\addtolength{\itemsep}{10pt}}{\endenumerate}
\newenvironment{widedescription}{\description\addtolength{\itemsep}{10pt}}{\enddescription}
\hypersetup{
colorlinks=true,
linkcolor=NTHU1,
filecolor=green, 
urlcolor=blue,
}
\beamertemplatenavigationsymbolsempty
\setbeamercolor{author in head/foot}{bg=white, fg=NTHU1}
\setbeamercolor{title in head/foot}{bg=white, fg=NTHU1}
\setbeamercolor{date in head/foot}{bg=white, fg=NTHU1}
\setbeamercolor{section in head/foot}{bg=white, fg=NTHU1}
\setbeamercolor{headline}{bg=white}
\setbeamertemplate{footline}{
    \leavevmode%
    \hbox{%
        \begin{beamercolorbox}[wd=.333333\paperwidth,ht=2.25ex,dp=1ex,center]{date in head/foot}%
            \usebeamerfont{date in head/foot}%\insertdate
        \end{beamercolorbox}%
        \begin{beamercolorbox}[wd=.333333\paperwidth,ht=2.25ex,dp=1ex,center]{title in head/foot}%
            \usebeamerfont{title in head/foot}\insertshorttitle
        \end{beamercolorbox}%
        \begin{beamercolorbox}[wd=.333333\paperwidth,ht=2.25ex,dp=1ex,right]{date in head/foot}%
            \usebeamerfont{date in head/foot}\insertframenumber{}\hspace*{2em}
        \end{beamercolorbox}}%
        \vskip0pt%
    }
    %% May need to script the introduction!!!!!!
%\setbeamercolor{page number in head/foot}{fg=NTHU1}
\setbeamertemplate{section in toc}[sections numbered]
\setbeamertemplate{subsection in toc}{\leavevmode\leftskip=3em\rlap{\hskip-1.75em\inserttocsectionnumber.\inserttocsubsectionnumber}
\inserttocsubsection\par}
\setbeamerfont{subsection in toc}{size=\footnotesize}
%\setbeamertemplate{headline}{%
  %\begin{beamercolorbox}[ht=5.5ex]{section in head/foot}
    %\vskip2pt\insertnavigation{0.33\paperwidth}\vskip2pt
  %\end{beamercolorbox}%
%}
\newenvironment{transitionframe}{\setbeamercolor{background canvas}{bg=white}\setbeamertemplate{footline}{} \begin{frame}[noframenumbering]}{\end{frame}}

\makeatletter
\let\@@magyar@captionfix\relax
\makeatother

\title[]{Conflict and Peace} % Change this regularly
\subtitle[]{Week 5: Impact of crime and conflict on individuals}
\author[Lee]{Seung-hun Lee }
\institute[]{Taipei School of Economics\\ National Tsing Hua University}

\date[]{October 3rd, 2025}

\begin{document}
\begin{transitionframe}
\titlepage
\end{transitionframe}

%%% Color slides for section headers: Use for colloquium version (The ones bewteen \iffals and \fi)

\section{Overview of the topic}
\begin{frame}
\frametitle{How do conflicts affect individuals?} 
Conflicts have very clear short-run implications
\begin{itemize}\setlength\itemsep{3pt}
\item Loss of life - both in terms of deaths and reduced healthy years
  \item Destruction of physical capital and infrastructures
\item Reduced economic activity and human capital accumulation
\end{itemize} \medskip
What is very worrisome: Negative effects can persist
\begin{itemize}\setlength\itemsep{3pt}
\item How do former combatants reintegrate?
\item Does it affect how we value future? ($\because$ Uncertainty about stable future)
\item Does it lock-in vulnerable individuals to crime / conflict specific capital?
\end{itemize}
\end{frame}

\begin{frame}
\frametitle{Some effects are already long-run!} 
Negative permanent health effects (Bundervoet et al 2009)
\begin{figure}
\includegraphics[keepaspectratio, width=0.4\textwidth]{fig1.png}
\caption{Progression of Civil war in Burundi}
\end{figure}
\end{frame}

\begin{frame}
\frametitle{Some effects are already long-run!} 
Negative permanent health effects (Bundervoet et al 2009)
\begin{figure}
\includegraphics[keepaspectratio, width=0.6\textwidth]{fig2.png}
\caption{Progression of Civil war in Burundi}
\end{figure}
\end{frame}

\begin{frame}
\frametitle{Conflicts can trigger other channels with long-run implications} 
Interrupting human capital accumulation
\begin{itemize}\setlength\itemsep{3pt}
\item Unable to acquire skills at schools
\item Could have implications in the long-run labor market performance
\item Potentially may acquire skills only applicable in limited (illicit) sectors
\end{itemize} \medskip
But studying this topic is very tricky
\begin{itemize}\setlength\itemsep{3pt}
\item Data is rare: Mostly developing countries + Dangerous to obtain them
\item Need to observe these individuals in the long-run
\item Even if acquired, endogeneity issues!
\begin{itemize}\setlength\itemsep{3pt}
  \item Ensuring that exposed individuals are comparable to the non-exposed are challenging
  \item Measurement error in being exposed to crime / conflict
\end{itemize}
\end{itemize} \medskip
\end{frame}

\begin{frame}
\frametitle{What we do today} 
Highlight novel channels that highlight effect of conflict in individuals
\begin{itemize}\setlength\itemsep{3pt}
\item Ability to make long-run investment $\Downarrow$
\begin{itemize}\setlength\itemsep{3pt}
\item Become more risk-averse
\item Uncertain environments (state presence, underlying policy changes)
\item Equally affect individuals and firm decisions
\end{itemize}
\item Accumulation of specific capital that does not apply to general skills
\begin{itemize}\setlength\itemsep{3pt}
\item Illegal industry specific capital: Difficult to look for other opportunities
\item Also potentially explain their criminal life path
\end{itemize}
\end{itemize} \medskip
We highlight potential channels that explain how effects of conflicts in individuals persist
\end{frame}

\begin{transitionframe}
\frametitle{Roadmap for today}
\tableofcontents
\end{transitionframe}

\section{Impact on individual and firm decision-making}
\subsection{Brown, Montalava, Thomas, Velásquez (2019): Impact of Violent Crime on Risk Aversion: Evidence from the Mexican Drug War}
\begin{transitionframe}
\begin{center}
  \begin{huge}
    \textcolor{NTHU1}{%
  Brown, Montalava, Thomas, Velásquez (2019): \\ Impact of Violent Crime on Risk Aversion: Evidence from the Mexican Drug War
    }
  \end{huge}
\end{center}
\end{transitionframe}

\begin{frame}
\frametitle{Q: How does exposure to crime affect attitudes towards risk} 
Attitudes towards risk influence key decisions over the course of one's life
\begin{itemize}\setlength\itemsep{3pt}
\item Willingness to take risks affect choices on investment, occupation, migration etc
\item Less known: How do life events change willingness to take risks?
\begin{itemize}\setlength\itemsep{3pt}
\item Some studies in Psychology dating 1990s (Carmil \& Breznitz 1991)
\end{itemize}
\item Causal analysis may be difficult
\begin{itemize}\setlength\itemsep{3pt}
\item Exposure to life events that can change attitudes may be selective to begin with
\item Thus, those exposed to such events may differ on many dimensions to those not exposed
\end{itemize}
\end{itemize}\medskip
This paper uses panel survey data and variation in crime exposure 
\begin{itemize}\setlength\itemsep{3pt}
\item Uses surveys conducted before and after surge of violence following War on Drugs
\item Leverages location, timing, and magnitude of violence as variation
\end{itemize}
\end{frame}

\begin{frame}
\frametitle{Violence in Mexico} 
Surge of violence in mid-2000s
\begin{itemize}\setlength\itemsep{3pt}
\item Homicide rate $<10$ per 100,000 before $\to$ $\simeq30$ after (currently 5th highest)
\item $\Uparrow$ police vs drug cartels, fights among splintered cartels since onset of War on Drugs
\item Splintered cartels spread throughout Mexico
\end{itemize} \medskip
How these events could affect attitudes towards risk
\begin{itemize}\setlength\itemsep{3pt}
\item Fear $\Uparrow$ $\to$ Optimism $\Downarrow$, Risk-taking $\Downarrow$ (Risk-averseness $\Uparrow$)
\item Economic conditions $\Downarrow$ $\to$ Risk-taking $\Downarrow$ (Risk-averseness $\Uparrow$)
\item Health $\Downarrow$ $\to$ Risk-taking $\Downarrow$ (Risk-averseness $\Uparrow$)
\item[$\to$] Goal is to distinguish these channels using data
\end{itemize}
\end{frame}

\begin{frame}
\frametitle{Data and empirical strategy} 
Data sources
\begin{itemize}\setlength\itemsep{3pt}
\item Individual panel surveys (Mexican Family Life Survey 2 and 3)
\begin{itemize}\setlength\itemsep{3pt}
  \item Socioeconomics and demographics
  \item Risk aversion measured with question on types returns from hypothetical gambles
  \item Attrition not correlated with conflict environment and other covariates
  \end{itemize}
\item Homicide statistics from government statistics office (INEGI)
\end{itemize} \medskip
Empirical strategy leverages violence across municipalities and time with panel regression
\[
y_{ijmt}=\beta_1 \text{Hom}_{jt}+\beta_2 X_{it-1}+\theta_i+\gamma_t + \lambda_m+e_{ijmt}
\]
\vspace{-10pt}
\begin{itemize}\setlength\itemsep{3pt}
  \item ($i,j,m,t$)=(indiv, origin, destination, year)
  \item $Y_{ijmt}=1$ if risk averse, $X_{it-1}$ measured year before to avoid reverse causality
  \item Attrition not correlated with $\text{Hom}_{jt}$ (Otherwise, $\beta_1$ is biased+-))
  \end{itemize}

\end{frame}

\begin{frame}
\frametitle{Leveraging pre vs post violent waves} 
\begin{figure}
\includegraphics[keepaspectratio,width=0.7\textwidth]{bmtv2019_f1.jpeg}
\end{figure}
\end{frame}

\begin{frame}
\frametitle{Eliciting risk from survey questions} 
\begin{figure}
\includegraphics[keepaspectratio,width=\textwidth]{bmtv2019_f2.jpeg}
\end{figure}
\end{frame}

\begin{frame}
\frametitle{Cross section vs panel regression discrepancy} 
\begin{columns}
\begin{column}{0.5\textwidth}
\begin{figure}
\includegraphics[keepaspectratio,width=\textwidth]{bmtv2019_t1.jpeg}
\end{figure}
\end{column}
\begin{column}{0.5\textwidth}
Why the discrepancy
\begin{itemize}\setlength\itemsep{3pt}
  \item MxFLS2: Less risk-averse indivs more likely to stay in high violence areas and agree to be surveyed
  \item Using panel: Observe same people across time - Able to use fixed effects to control unobservable characteristics
\end{itemize}
\end{column}
\end{columns}

\end{frame}

\begin{frame}
\frametitle{Fear seems to be the most relevant mechanism} 
\begin{figure}
\includegraphics[keepaspectratio,width=\textwidth]{bmtv2019_t2.jpeg}
\end{figure}
\end{frame}

\begin{frame}
\frametitle{Takeaways and discussions} 
Violence can affect decision-making by affecting risk attitudes
\begin{itemize}\setlength\itemsep{3pt}
\item Violence exposure $\to$ Increased risk-averseness (decreased risk-taking)
\item Driven by increased sense of fear following violence
\item Compares with conventional cross-section estimates: Show that they can be biased
\end{itemize}\medskip
What can we take away from here?
\begin{itemize}\setlength\itemsep{3pt}
\item Consequence of violence: Less willingness to engage in risky but beneficial endeavors
\begin{itemize}\setlength\itemsep{3pt}
\item Could lead to lot of missed opportunities for development
\item Which implies less wealth accumulation and income
\end{itemize}
\item Another channel in which violence hurts development
\begin{itemize}\setlength\itemsep{3pt}
\item Violence need not destroy capital directly to impact growth and development
\end{itemize}
\end{itemize}
\end{frame}

\subsection{Blair, Christensen, Wirtschafter (2022): How Does Armed Conflict Shape Investment? Evidence from the Mining Sector}
\begin{transitionframe}
\begin{center}
  \begin{huge}
    \textcolor{NTHU1}{%
  Blair, Christensen, Wirtschafter (2022): \\ How Does Armed Conflict Shape Investment? Evidence from the Mining Sector
    }
  \end{huge}
\end{center}
\end{transitionframe}

\begin{frame}
\frametitle{Q: How do conflicts affect firms' investment decisions?} 
Uncertainty can disincentivize investment
\begin{itemize}\setlength\itemsep{3pt}
\item Armed conflicts increase uncertainties over future economic returns
\item Many measures taken to reduce armed conflict
\begin{itemize}
\item e.g. Monopolize violence (only the state can legally control military/police etc)
\end{itemize}
\item But do armed conflicts have unilateral effects on investment incentives? 
\end{itemize}\medskip
This paper tells us that armed conflicts do have divergent effects on investments
\begin{itemize}\setlength\itemsep{3pt}
\item Even within the same country, distance from conflict may affect changes to incentives
\item Destruction of production, state absence, uncertain policies apply differently
\begin{itemize}
  \item Local threats (Destruction, state absence) vs Universal factors (uncertainty) 
\end{itemize}
\item Reviews findings from previously publish papers 
\end{itemize}\medskip
\end{frame}

\begin{frame}
\frametitle{Theoretical framework of conflict and investments} 
Firms concerned about political instability to the extent that it constrains operations
\begin{itemize}\setlength\itemsep{3pt}
\item Factors to consider: Production closure, state absence, uncertainty, reputation
 \begin{itemize}\setlength\itemsep{3pt}
  \item Production closure (-)
  \item State absence: Lax regulations, difficult to tax (+), less protection of property (-)
  \item Uncertain policies, government action and income horizon (-)
  \item Reputation for working with warring government (-)
\end{itemize} \medskip
\item Exposure to political risk is heterogenous (by proximity)
\begin{itemize}\setlength\itemsep{3pt}
  \item Conflict site: Destruction of production sites $\to$ Investment $\Downarrow$
  \item Buffer zone: Ungoverned and difficult to tax $\to$ Investment?
  \item Country with Conflict: Still governed normally $\to$ Investment?
  \end{itemize}
\end{itemize}\medskip
\end{frame}

\begin{frame}
\frametitle{Theoretical framework of conflict and investments} 
\begin{figure}
\includegraphics[keepaspectratio,width=0.7\textwidth]{bcw2022_f1.jpeg}
\end{figure}
\end{frame}

\begin{frame}
\frametitle{Theoretical framework of conflict and investments} 
\begin{figure}
\includegraphics[keepaspectratio,width=0.9\textwidth]{bcw2022_t1.jpeg}
\end{figure}
\end{frame}

\begin{frame}
\frametitle{Data and empirical strategy} 
Data sources
\begin{itemize}\setlength\itemsep{3pt}
\item Investment: Mining investment from SNL Metals and Mining dataset
\begin{itemize}
\item Includes data on investments made, locations from 1997-2014
\end{itemize}
\item Conflict data from UCDP dataset
\end{itemize}\medskip
Research design leverages variation in distance from conflict site and investment
\[
y_{ict}=\alpha_{ic}+\delta_{it} + \theta C_{ct} + \sum_k \kappa^k D^k_{ict}+e_{ict}
\]\vspace{-15pt}
\begin{itemize}\setlength\itemsep{3pt}  
\item firm-country, firm-time fixed effects
\item $C_{ct}$ capture whether a conflict happened at country $c$ in year $t$
\item $k$ is a distance bandwidth
\item $ D^k_{ict}$ captures whether firm $i$ in $c$ is exposed to conflict at year $t$ in bandwidth $k$ 
\end{itemize}
\end{frame}


\begin{frame}
\frametitle{Effects differ by proximity} 
\begin{figure}
\includegraphics[keepaspectratio,width=0.48\textwidth]{bcw2022_f2.jpeg}
\end{figure}
\end{frame}

\begin{frame}
\frametitle{Yet these heterogeneities are not captured at aggregate level} 
\begin{figure}
\includegraphics[keepaspectratio,width=0.6\textwidth]{bcw2022_t2.jpeg}
\end{figure}
\end{frame}

\begin{frame}
\frametitle{Takeaways and discussions} 
Conflicts do not have unilateral influence on investment outcomes
\begin{itemize}\setlength\itemsep{3pt}
\item Different factors at work depending on proximity to conflict 
\item Can take advantage of lax rules due to state absence? they may increase investments
\item Otherwise, conflicts can discourage investment
\end{itemize}\medskip
Possible extensions
\begin{itemize}\setlength\itemsep{3pt}
\item Effects of state absence on investments (later: Sanchez de la Sierra (2020))
\begin{itemize}\setlength\itemsep{3pt}
\item Non-state actors can take over (call them bandits)
\item If protecting business activity is in bandits' interests, economic activity can continue
\item Otherwise, they have no interest to protect
\end{itemize} \medskip
\item This difference could appear in dimensions other than proximity
\begin{itemize}\setlength\itemsep{3pt}
\item Difference by sector of industry, for instance
\end{itemize}
\end{itemize}\medskip
\end{frame}


\subsection{Mirenda, Mocetti, Rizzica (2022): The Economic Effects of Mafia: Firm Level Evidence[Student presentation]}
\begin{transitionframe}
\begin{center}
  \begin{huge}
    \textcolor{NTHU1}{%
Mirenda, Mocetti, Rizzica (2022): \\ The Economic Effects of Mafia: Firm Level Evidence [Student presentation]
    }
  \end{huge}
\end{center}
\end{transitionframe}

 

\section{Effects on lifepath of individuals}
\subsection{Humphreys, Weinstein (2007): Demobilization and Reintegration}
\begin{transitionframe}
\begin{center}
  \begin{huge}
    \textcolor{NTHU1}{%
  Humphreys, Weinstein (2007):\\ Demobilization and Reintegration
    }
  \end{huge}
\end{center}
\end{transitionframe}

\begin{frame}
\frametitle{Q: What determines successful demobilization and reintegration?} 
Many efforts to end protracted conflicts and bring normalcy afterwards
\begin{itemize}\setlength\itemsep{3pt}
\item Peacekeeping operations (PKOs) led by the United Nations
\begin{itemize}\setlength\itemsep{3pt}
\item Reform institutions and rebuild infrastructures to reduce incentives for further violence
\end{itemize}
\item Disarmament, demobilization, and reintegration (DDR) is key part of this process
 \begin{itemize}\setlength\itemsep{3pt}
\item  El Salvador, Cambodia, Sierra Leone, Tajikistan etc
\end{itemize}
\item Not enough micro-level evidence to assess the success
\end{itemize}\medskip
This paper uses micro-level evidence to identify factors contributing DDR programs
 \begin{itemize}\setlength\itemsep{3pt}
\item Survey of combatants from warring factions in Sierra Leone
\item Variation in the degree of reintegration across former combatants
\item Identifies factors - personal characteristics, and contexts - that explains DDR effects
\end{itemize}
\end{frame}

\begin{frame}
\frametitle{Setting: Post-civil war in Sierra Leone (SL)} 
History of post-war progress
 \begin{itemize}\setlength\itemsep{3pt}
\item War declared over in January 2002 by SL government (after 11 years!)
\item 5 warring entities involved - War ended with capture of leader of a major faction
 \end{itemize}\medskip
 DDR process 
 \begin{itemize}\setlength\itemsep{3pt}
\item Assemble at reception center $\to$ Disarm $\to$ receive allowances, necessities $\to$ civilian
\item Receive job training (equally across 5 factions, uncorrelated with political affiliations)
\end{itemize}\medskip

Hypotheses on successes of DDR
 \begin{itemize}\setlength\itemsep{3pt}
\item Success depends on dissolving combatant networks, improving earnings, fostering trust on democratic process, and reconciling with families and friends
\item Those not happy with earnings potential + distrust other groups less likely to integrate
\item Conversely, DDR could be most pronounced on high dissatisfaction and distrust indivs
\end{itemize}
\end{frame}

\begin{frame}
\frametitle{Data and empirical strategy} 
Outcome variables from surveying 1000+ combatants across warring groups
 \begin{itemize}\setlength\itemsep{3pt}
\item Delinked: Break ties with the faction? (time/business relation/seeking help)
\item Employed: Reintegrate to workforce
\item Democratic: Confident in democracy? (elect/talk to officials vs fight/traditional leader)
\item Accepted: Any problems in gaining acceptance from family and friends?
\end{itemize}\medskip
Key independent variables of integration
 \begin{itemize}\setlength\itemsep{3pt}
\item Distrust: Believe in others to comply to peace agreement?
\item Dissatisfaction: Faction benefitted or worse-off after agreement?
\end{itemize}\medskip
The paper estimate correlates across various characteristics and above outcomes
 \begin{itemize}\setlength\itemsep{3pt}
\item Probit with dummies for factions applied
\item Not necessarily causal: Not a randomized procedure in allocating one to DDR
 \begin{itemize}\setlength\itemsep{3pt}
\item Refer to Blattman and Berman's works we covered few weeks ago
\end{itemize}
\end{itemize}
\end{frame}

\begin{frame}
\frametitle{How are the combatants doing in these outcomes?}
\begin{figure}
  \includegraphics[keepaspectratio,width=0.9\textwidth]{mh2007_t1.jpeg}
\end{figure} 
\end{frame}

\begin{frame}
\frametitle{Who are more or less likely to reintegrate to begin with?}
\begin{columns}
\begin{column}{0.5\textwidth}
  \begin{figure}
  \includegraphics[keepaspectratio,width=\textwidth]{mh2007_t3.jpeg}
\end{figure} 
\end{column}

\begin{column}{0.5\textwidth}
\begin{itemize}\setlength\itemsep{5pt}
\item \textcolor{blue}{Distrust and dissatisfaction} only explains a small part of reintegration
\item Some work in opposite directions (\textcolor{red}{political support for factions})
\item \textcolor{green}{High-ranked, wealthier, and educated} ones less likely to reintegrate (on some dimensions only)
\end{itemize}
\end{column}
\end{columns}

\end{frame}

\begin{frame}
\frametitle{No systematic evidence of the effect of DDR on reintegration}
\begin{figure}
  \includegraphics[keepaspectratio,width=0.9\textwidth]{mh2007_t2.jpeg}
\end{figure} 
\end{frame}

\begin{frame}
\frametitle{Takeaways and discussions}
The observed correlation between DDR and reintegration is weak 
\begin{itemize}\setlength\itemsep{3pt}
\item Dissatisfaction and distrust only explains partially
\item Some personal characteristics are more robustly correlated
\item Participation at DDR does not contribute to reintegration\
\end{itemize}\medskip
So is DDR not effective at all? Not so fast...
\begin{itemize}\setlength\itemsep{3pt}
\item Differential factors within DDR could have different effects
\begin{itemize}\setlength\itemsep{3pt}
  \item Difference in providing short-run money vs long-run opportunities vs cognitive care
  \item See Blattman, Jamison, Sheridan (2017) - Companion of Blattman, Annan (2016) 
\end{itemize}
\item Could be different depending on how conflict ended
\begin{itemize}\setlength\itemsep{3pt}
  \item This one ended with one side losing
  \item Not all conflicts end this way (most end without clear winners)
\end{itemize}
\end{itemize}
\end{frame}

\subsection{Sviatchi (2022): Making a Narco: Childhood Exposure to Illegal Labor Markets and Criminal Life Paths}
\begin{transitionframe}
\begin{center}
  \begin{huge}
    \textcolor{NTHU1}{%
  Sviatchi (2022):\\ Making a Narco: Childhood Exposure to Illegal Labor Markets and Criminal Life Paths
    }
  \end{huge}
\end{center}
\end{transitionframe}

\begin{frame}
\frametitle{Q: Does early exposure to criminal life affect future life paths?} 
Illegal markets operated by organized criminal groups involve so many individuals
\begin{itemize}\setlength\itemsep{3pt}
\item The `Kingpins': Pablo Escobar, Arturo Beltran Leyva, Al Capone ..
\item But there are also the `footsoldiers': Mostly from marginalized backgrounds
\item So How did they enter the criminal life?
\begin{itemize}
\item Spot benefit-cost calculus? Persistent factors that lock them into criminal life?
\end{itemize}
\end{itemize}\medskip
This paper: Early exposure $\to$ formation of very specific human capital $\to$ lock-in
\begin{itemize}\setlength\itemsep{3pt}
\item Exogenous shocks to illegal markets in Peru from US intervention in Colombia
\item Return to illegal act $\Uparrow$$\to$ Child labor in coca farming $\Uparrow$
\item Accumulate human capital specific to illegal activities $\to$ Incarceration $\Uparrow$
\item Early interventions (cash transfers for education) may offset these effective
\end{itemize}\medskip
\end{frame}

\begin{frame}
\frametitle{Illegal drug markets in Latin America} 
Coca Production Colombia, Bolivia, Peru: Only places agroecologically suitable
\begin{itemize}\setlength\itemsep{3pt}
\item Change of drug policy in one can spill over, but to one of these countries only!
\item As US shifted anti-drug efforts to Colombia, production shifted to Peru (early 2000s)
\item Aerial spraying at Colombia hiked prices $\to$ Coca production rises in Peru
\item Largest producer of cocaine by 2012 
\end{itemize}\medskip
Which individuals are more likely to be `locked-in' to this sector?
\begin{itemize}\setlength\itemsep{3pt}
\item Labor market situation matters when deciding to continue education or work
\item Age 11-14: End of primary school $\to$ secondary education vs labor market
\item Coca price hike has effect on this group: Induce to switch off from education
\end{itemize}\medskip
\end{frame}

\begin{frame}
\frametitle{Data sources} 
Data sources
\begin{itemize}\setlength\itemsep{3pt}
\item Coca and other agricultural data: Satellite imagery + Agricultural census
\item Coca price: UN Office on Drugs and Crime
\item 2 Peruvian household surveys: Educational attainment and labor outcomes
\item School and incarceration data obtained from administrative data
\end{itemize}\medskip
Empirical strategy leverages price shocks, timing of exposure, and location
\begin{itemize}\setlength\itemsep{3pt}
\item Labor market effects for indiv $i$ at area $d$ at time $t$ in age  bracket $k$: $Y_{idt} = \beta_0(PC_t\times \text{Coca}_d) + \sum_k \beta_k \text{Age}_k \times (PC_t\times \text{Coca}_d) +\alpha_d + \phi_t + \sigma_r t + X_{idt}+e_{idt}$ \medskip
\item Long-run effects (incarceration and crime) for each birth cohort $c$  $Y_{dc}=\beta\text{Priceshock}\times\text{Age11-14}_{dc}+\sum_k \beta_k\text{Priceshock}\times\text{Age $k$}_{dc} + \alpha_d + \delta_c + \sigma_d c + e_{cd}$ \medskip
\item Why? Incarceration data at cohort level (not individuals)
\end{itemize}\medskip
\end{frame}

\begin{frame}
\frametitle{Geographic and Temporal variation} 
\begin{columns}
\begin{column}{0.5\textwidth}
\begin{figure}
\includegraphics[keepaspectratio, width=\textwidth]{s2022_f1.jpeg}
\caption{Coca production over time}
\end{figure}
\end{column}
\begin{column}{0.5\textwidth}
  \begin{figure}
\includegraphics[keepaspectratio, width=\textwidth]{s2022_f2.jpeg}
\caption{Coca suitability (1994) and production intensity (2000s)}
\end{figure}
\end{column}
\end{columns}
\end{frame}

\begin{frame}
  \frametitle{Price shock increase illegal labor participation, esp for 11-14 yr olds}
    \begin{figure}
\includegraphics[keepaspectratio, width=0.8\textwidth]{s2022_t1.jpeg}
\end{figure}
\end{frame}

\begin{frame}
  \frametitle{This exposure locks them into criminal industry-specific capital}
    \begin{figure}
\includegraphics[keepaspectratio, width=0.7\textwidth]{s2022_t2.jpeg}
\end{figure}
\end{frame}

\begin{frame}
  \frametitle{Timely educational intervention offsets lock-in effects}
    \begin{figure}
\includegraphics[keepaspectratio, width=0.55\textwidth]{s2022_t3.jpeg}
\end{figure}
\end{frame}

\begin{frame}
\frametitle{Takeaways and Discussion} 
Exposure to illegal markets and crime locks children into criminal life path
\begin{itemize}\setlength\itemsep{3pt}
\item Drop out of school: Cannot accumulate formal human capital
\item More likely to be incarcerated later
\item Intervention offsetting incentives to join labor early do mitigate some effects
\end{itemize}\medskip
Discussion 
\begin{itemize}\setlength\itemsep{3pt}
\item Hints that environment where children grows matters
\begin{itemize}
\item Active research in economics, education, psychology (Moving to Opportunity etc)
\end{itemize}
\item Exposure at critical timing can tilt incentives, leading to long-run effects 
\begin{itemize}
\item Exposure to adverse events (crime, natural disasters) in utero
\end{itemize}
\end{itemize}\medskip
\end{frame}


\subsection{Couttenier, Petrencu, Rohner, Thoenig (2019): The Violent Legacy of Conflict: Evidence on Asymlum Seekers, Crime, and PublicPolicy in Switzerland [Student Presentation]}
\begin{transitionframe}
\begin{center}
  \begin{huge}
    \textcolor{NTHU1}{%
Couttenier, Petrencu, Rohner, Thoenig (2019): \\ The Violent Legacy of Conflict: Evidence on Asymlum Seekers, Crime, and Public Policy in Switzerland [Student Presentation]
    }
  \end{huge}
\end{center}
\end{transitionframe}

\begin{frame}
\frametitle{Takeaways and remaining questions} 
Conflicts affect individuals through many channel
\begin{itemize}\setlength\itemsep{3pt}
\item Direction destruction of physical capital and life opportunities
\begin{itemize}
\item Not directly covered this week, but indirectly hinted throughout semester
\end{itemize}
\item Affect perspectives on future outlook through changing attitudes towards risk
\item May lock-in individuals exposed to crime, especially at critical timing
\end{itemize} \medskip
Remaining questions
\begin{itemize}\setlength\itemsep{3pt}
\item What is the effective rehabilitation and reintegration methods 
\begin{itemize}
\item Not just conflicts, but applies to criminals (Law and economics, labor economics)
\end{itemize}
\item Any institutional setup that mitigate or aggravates impact on individuals? (Next week)
\end{itemize}
\end{frame}


%%%%%%%%%%%
\end{document}


